\section{USAMO 2010}
        \begin{enumerate}
        	\item (\textit{USAMO 2010}) Let $AXYZB$ be a convex pentagon inscribed in a semicircle of diameter $AB$. Denote by $P$, $Q$, $R$, $S$ the feet of the perpendiculars from $Y$ onto lines $AX$, $BX$, $AZ$, $BZ$, respectively.  Prove that the acute angle formed by lines $PQ$ and $RS$ is half the size of $\angle XOZ$, where $O$ is the midpoint of segment $AB$.


 

	\item (\textit{USAMO 2010}) There are $n$ students standing in a circle, one behind the other. The students have heights $h_1<h_2<\dots <h_n$. If a student with height $h_k$ is standing directly behind a student with height $h_{k-2}$ or less, the two students are permitted to switch places. Prove that it is not possible to make more than $\binom{n}{3}$ such switches before reaching a position in which no further switches are possible.


 

	\item (\textit{USAMO 2010}) The 2010 positive numbers $a_1, a_2, \ldots , a_{2010}$ satisfy the inequality $a_ia_j\leq i+j$ for all distinct indices $i, j$. Determine, with proof, the largest possible value of the product $a_1a_2\ldots a_{2010}$.


 

	\item (\textit{USAMO 2010}) Let $ABC$ be a triangle with $\angle A=90^{\circ}$. Points $D$ and $E$ lie on sides $AC$ and $AB$, respectively, such that $\angle ABD=\angle DBC$ and $\angle ACE=\angle ECB$. Segments $BD$ and $CE$ meet at $I$. Determine whether or not it is possible for segments $AB$, $AC$, $BI$, $ID$, $CI$, $IE$ to all have integer side lengths.


 

	\item (\textit{USAMO 2010}) Let $q = \dfrac{3p-5}{2}$ where $p$ is an odd prime, and let


 

 \[S_q =    \frac{1}{2\cdot 3 \cdot 4} +   \frac{1}{5\cdot 6 \cdot 7} + \cdots +   \frac{1}{q\cdot (q+1) \cdot (q+2)}.\]


 


 Prove that if $\dfrac{1}{p}-2S_q = \dfrac{m}{n}$ for relatively prime integers
$m$ and $n$, then $m-n$ is divisible by $p$.


 

	\item (\textit{USAMO 2010}) A blackboard contains 68 pairs of nonzero integers.  Suppose that for each positive integer $k$ at most one of the pairs $(k, k)$ and $(-k, -k)$ is written on the blackboard.  A student erases some of the 136 integers, subject to the condition that no two erased integers may add to 0.  The student then scores one point for each of the 68 pairs in which at least one integer is erased.  Determine, with proof, the largest number $N$ of points that the student can guarantee to score regardless of which 68 pairs have been written on the board.


 

\end{enumerate}